\section{Lean verification}\label{sec:lean}

\subsection{Repository and build}

The anonymous repository is
\begin{center}
  \texttt{a-stringflow/stringflow-proof}
\end{center}
with the Lean project in the \texttt{lean/} subdirectory.  The verified
toolchain is documented in the repository
\cite{stringflowproof}.
\begin{verbatim}
Lean 4 : v4.33.0-rc2
Mathlib: locked by lakefile
build  : lake build CycleBridge
         lake build DwdbDiv
\end{verbatim}

\subsection{File map}

\begin{center}
\begin{tabular}{ll}
\toprule
File & Role \\
\midrule
\texttt{FinitePrefix.lean} & explicit 17-state expansion \\
\texttt{UnifiedCoreBridge.lean} & candidate and segment bridges \\
\texttt{UnifiedCoreAudit.lean} & unified-core audit \\
\texttt{RealOrbitCharge.lean} & cyclic rise block decomposition and charge balance \\
\texttt{amiya.lean} & run budget bounds and tail rank threshold \\
\texttt{trinity.lean} & trinity block and failure window route \\
\texttt{Closure.lean} & four-term additive identity and all-below-budget crush \\
\texttt{DwdbDiv.lean} & final divergence assembly \\
\texttt{FinalTheorem.lean} & public divergence statement \texttt{five\_x\_plus\_one\_diverges\_at\_7} \\
\bottomrule
\end{tabular}
\end{center}

\subsection{Repository layout}

The main files under \texttt{lean/} are:
\begin{verbatim}
WordWindow.lean          word, orbit, validity, molecule
Pmi.lean                 PMI and PMI-B algebra
FinitePrefix.lean        explicit finite base
UnifiedCoreBridge.lean   candidate and segment bridges
RealOrbitCharge.lean     cyclic block decomposition and charge balance
amiya.lean               run budget and rank threshold theorems
trinity.lean             trinity blocks and failure window routing
Closure.lean             four-term additive identity and crush theorems
DwdbDiv.lean             final divergence assembly
FinalStatement.lean      public divergence definition
FinalTheorem.lean        final assembled theorem
AxiomAudit.lean          axiom audit
\end{verbatim}

\subsection{Statement inventory by file}

\begin{center}
\scriptsize
\setlength{\tabcolsep}{4pt}
\begin{tabular}{@{}p{0.26\textwidth}p{0.58\textwidth}p{0.12\textwidth}@{}}
\toprule
File & Key statements & Status \\
\midrule
\texttt{FinitePrefix.lean} &
  \texttt{fullOrbitIter\_\allowbreak prefix\_\allowbreak expand},
  \texttt{fullOrbit\_\allowbreak prefix\_\allowbreak step\_\allowbreak
  weights},
  \texttt{fullOrbit\_\allowbreak first\_\allowbreak t\_\allowbreak
  ge3\_\allowbreak is\_\allowbreak exactly\_\allowbreak 3},
  \texttt{first\_\allowbreak big\_\allowbreak step\_\allowbreak unique}
  & compiled \\
\texttt{UnifiedCoreBridge.lean} &
  \texttt{candidateX},
  \texttt{reset\_\allowbreak head\_\allowbreak predecessor},
  \texttt{candidateX\_\allowbreak of\_\allowbreak reset\_\allowbreak
  and\_\allowbreak terminal},
  \texttt{first\_\allowbreak block\_\allowbreak terminal\_\allowbreak eq},
  \texttt{d1\_\allowbreak exclusion},
  \texttt{d2\_\allowbreak exclusion},
  \texttt{d3\_\allowbreak family\_\allowbreak big\_\allowbreak
  weight\_\allowbreak excluded},
  \texttt{dge4\_\allowbreak e2\_\allowbreak a\_\allowbreak ge1\_\allowbreak
  excluded} & compiled \\
\texttt{UnifiedCoreAudit.lean} &
  \texttt{All36\_\allowbreak 20PremisesNoHge},
  \texttt{blockB\_\allowbreak bound\_\allowbreak of\_\allowbreak
  no\_\allowbreak hge},
  \texttt{rj0\_\allowbreak ge\_\allowbreak of\_\allowbreak
  size\_\allowbreak conditions\_\allowbreak no\_\allowbreak hge},
  \texttt{terminal\_\allowbreak bound\_\allowbreak iff\_\allowbreak
  yStar\_\allowbreak ne\_\allowbreak zero},
  \texttt{unified\_\allowbreak core\_\allowbreak final\_\allowbreak
  no\_\allowbreak hge} & compiled \\
\texttt{CycleBridge.lean} &
  \texttt{cycleWord},
  \texttt{CycleQb8Input},
  \texttt{rise\_\allowbreak block\_\allowbreak balance},
  \texttt{failure\_\allowbreak window\_\allowbreak contradicts\_\allowbreak
  window\_\allowbreak corrected},
  \texttt{windowBoundToNoCycle\_\allowbreak of\_\allowbreak
  failureWindowExistence} & compiled \\
\texttt{DwdbDiv.lean} &
  \texttt{dwdbDivFinalAssembly},
  \texttt{five\_x\_plus\_one\_diverges\_at\_7\_of\_failure\_window} & compiled \\
\texttt{FinalTheorem.lean} &
  \texttt{five\_x\_plus\_one\_diverges\_at\_7} & compiled \\
\texttt{FinalStatement.lean} &
  \texttt{fiveXPlusOneStep},
  \texttt{fiveXPlusOneOrbit},
  \texttt{IsUnboundedOrbit},
  \texttt{OrbitCycle} & compiled \\
\bottomrule
\end{tabular}
\end{center}

The inventory mirrors the theorem index in the supplementary manual.
Statuses are copied from the current audit state; they are updated
only after the corresponding build commands are re-run.

\subsection{Top-level definitions}

The four objects used in the main theorem are defined in
\texttt{FinalStatement.lean}:
\begin{verbatim}
def fiveXPlusOneStep (x : Nat) : Nat :=
  (5 * x + 1) / 2 ^ twoValuation (5 * x + 1)

def fiveXPlusOneOrbit (x : Nat) : Nat -> Nat
  | 0 => x
  | n + 1 => fiveXPlusOneStep (fiveXPlusOneOrbit x n)

def IsUnboundedOrbit (x : Nat) : Prop :=
  forall B : Nat, exists n : Nat, B <= fiveXPlusOneOrbit x n

def OrbitCycle (x : Nat) : Prop :=
  exists n m : Nat, m < n /\
    fiveXPlusOneOrbit x m = fiveXPlusOneOrbit x n
\end{verbatim}

The bridge theorem is
\begin{verbatim}
unbounded_of_no_cycle (x : Nat) (h : not OrbitCycle x) :
  IsUnboundedOrbit x
\end{verbatim}

\subsection{How to read the Lean code}

The Lean files are best read in dependency order.
\begin{enumerate}[leftmargin=*]
\item \texttt{FinitePrefix.lean}: explicit orbit values and weights;
\item \texttt{UnifiedCoreBridge.lean}: candidate parameterisation and
  segment bridges;
\item \texttt{UnifiedCoreAudit.lean}: CRT, terminal residue, and the
  final pending inequality;
\item \texttt{CycleBridge.lean}: corrected window bound and the
  no-cycle bridge;
\item \texttt{DwdbDiv.lean}: final divergence assembly.
\end{enumerate}
Each file states its dependencies in its header and in the
supplementary manual.

\subsection{Bridge interface map}

The bridge layer is a sequence of interfaces.  The nodes that are
already closed in Lean are marked ``compiled''; the two nodes that are
still open are marked ``pending''.  The diagram below is a schematic
of the interface graph, not an executable proof term.

\begin{verbatim}
FinitePrefix.lean (compiled)
  -> UnifiedCoreBridge.lean (compiled)
  -> UnifiedCoreAudit.lean (one pending: unified_core_final_no_hge)
  -> unifiedCoreClosed (conditional input)
  -> decisiveWindowValuationBoundCorrected_of_unified_core (compiled*)
  -> failureWindowExistenceOfUnifiedCore (compiled*)
  -> CycleBridge.windowBoundToNoCycle_of_failureWindowExistence
       (compiled*, conditional)
  -> DwdbDiv.dwdbDivFinalAssembly (compiled*, conditional)
  -> FinalStatement.five_x_plus_one_diverges_at_7 (pending)
\end{verbatim}

The starred nodes compile as conditional interfaces: each consumes the
preceding open input and does not assert that the input is already
proved.  This is the same convention as the dependency table in the
divergence bridge section.

\begin{center}
\small
\begin{tabular}{@{}p{0.44\textwidth}p{0.30\textwidth}p{0.20\textwidth}@{}}
\toprule
Interface & Consumes & Status \\
\midrule
\texttt{decisiveWindowValuation\allowbreak
BoundCorrected\_\allowbreak of\_\allowbreak unified\_\allowbreak core}
& \texttt{unifiedCoreClosed} & compiled* \\
\texttt{failureWindowExistenceOfUnifiedCore} &
  \texttt{unifiedCoreClosed} & compiled* \\
\texttt{windowBoundToNoCycle\_\allowbreak of\_\allowbreak
failureWindowExistence} &
  \texttt{failureWindowExistence} & compiled* \\
\texttt{dwdbDivFinalAssembly} &
  \texttt{windowBoundToNoCycle} & compiled* \\
\texttt{five\_\allowbreak x\_\allowbreak plus\_\allowbreak one\_\allowbreak
diverges\_\allowbreak at\_\allowbreak 7} &
  \texttt{dwdbDivFinalAssembly} & pending \\
\bottomrule
\end{tabular}
\end{center}

The map makes the remaining work explicit: closing
\texttt{unified\_\allowbreak core\_\allowbreak final\_\allowbreak
no\_\allowbreak hge} provides \texttt{unifiedCoreClosed}, which then
flows through the compiled conditional interfaces to the final
assembly target.

\subsection{Verification targets}

\begin{center}
\small
\begin{tabular}{@{}p{0.32\textwidth}p{0.34\textwidth}p{0.28\textwidth}@{}}
\toprule
Target & Command & Status \\
\midrule
\texttt{CycleBridge} & \texttt{lake build CycleBridge} & compiled \\
\texttt{DwdbDiv} & \texttt{lake build DwdbDiv} & compiled \\
\texttt{UnifiedCoreAudit} & \texttt{lake build UnifiedCoreAudit} &
contains \texttt{sorry} \\
\texttt{FinalStatement} & \texttt{lake build FinalStatement} &
assembly target \\
\bottomrule
\end{tabular}
\end{center}

\subsection{Acceptance criteria}

\begin{enumerate}[leftmargin=*]
\item No \texttt{sorry} in any theorem that the main chain consumes.
\item No custom \texttt{axiom}; the allowed axioms are only
  \texttt{propext}, \texttt{Classical.choice}, and \texttt{Quot.sound}.
\item The main chain contains no \texttt{native\_decide} trust; finite bases
  are expanded explicitly.
\item The final theorem is
  \begin{verbatim}
FinalStatement.five_x_plus_one_diverges_at_7 :
  IsUnboundedOrbit 7
  \end{verbatim}
\end{enumerate}

\subsection{Axiom audit}

The file \texttt{AxiomAudit.lean} prints the axioms of the main bridge
theorems.  For the corrected window-bridge theorems, the audit returns
only
\begin{verbatim}
[propext, Classical.choice, Quot.sound]
\end{verbatim}

\begin{center}
\small
\begin{tabular}{@{}p{0.55\textwidth}p{0.38\textwidth}@{}}
\toprule
Theorem & Axioms \\
\midrule
\texttt{no\_cycle\_of\_window\_bound} &
\texttt{[propext, Quot.sound]} \\
\texttt{failure\_window\_contradicts\_window\_corrected} &
\texttt{[propext, Quot.sound]} \\
\texttt{five\_x\_plus\_one\_diverges\_at\_7\_of\_window\_bound} &
\texttt{[propext, Classical.choice, Quot.sound]} \\
\texttt{DwdbDiv.dwdbDivFinalAssembly} &
\texttt{[propext, Classical.choice, Quot.sound]} \\
\bottomrule
\end{tabular}
\end{center}

The audit is re-run whenever the bridge is changed.  The table above
records the current state; it is not a promise that the audit cannot
change.

\subsection{Axiom audit methodology}

\texttt{AxiomAudit.lean} runs \texttt{\#print axioms} on each listed
bridge theorem and compares the printed axiom set with the allowlist
\texttt{\{propext, Classical.choice, Quot.sound\}}.  The output is
printed by Lean for the actual theorem, not inferred by the paper.  The
check is repeated whenever the dependencies change.

\subsection{Verification methodology: worked example}

The audit of one bridge theorem is a literal sequence of commands.
First the target is built:
\begin{verbatim}
lake build CycleBridge
\end{verbatim}
Then the axiom set of the assembled bridge is printed:
\begin{verbatim}
#print axioms CycleBridge.windowBoundToNoCycle
\end{verbatim}
The expected output is exactly
\begin{verbatim}
[propext, Classical.choice, Quot.sound]
\end{verbatim}
or the smaller set
\begin{verbatim}
[propext, Quot.sound]
\end{verbatim}
for theorems that do not use choice.  Any extra entry, in particular a
custom \texttt{axiom}, fails the audit.

The \texttt{sorry} check is a search over the audited files.  The
current state of the unified-core audit is exactly one occurrence:
\begin{verbatim}
UnifiedCoreAudit.lean: sorry
\end{verbatim}
at \texttt{unified\_core\_final\_no\_hge}.  The \texttt{native\_decide}
check searches the main-chain files for the token
\texttt{native\_decide}; the current count is zero.  Both searches are
run before a status table is updated.

\begin{center}
\small
\begin{tabular}{@{}p{0.40\textwidth}p{0.52\textwidth}@{}}
\toprule
Audit command & Expected current result \\
\midrule
\texttt{lake build CycleBridge} & success \\
\texttt{lake build DwdbDiv} & success \\
\texttt{\#print axioms} on bridge theorems &
  only \texttt{propext}, \texttt{Classical.choice}, \texttt{Quot.sound} \\
search for \texttt{sorry} in \texttt{UnifiedCoreAudit.lean} &
  exactly one occurrence \\
search for \texttt{native\_decide} in the main chain &
  zero occurrences \\
\bottomrule
\end{tabular}
\end{center}

This table is re-run whenever the repository changes.  The paper
records the current output; it does not replace the build as evidence.

\subsection{Continuous verification}

The repository is configured so that a clean checkout runs
\begin{verbatim}
lake build CycleBridge
lake build DwdbDiv
\end{verbatim}
as part of the verification workflow.  The anonymous repository and its
CI status are recorded in \cite{stringflowproof}.

\subsection{What a successful build establishes}

A successful \texttt{lake build} checks that every theorem in the
target files type-checks in the pinned environment.  It does not run
orbit simulations, and it does not itself decide whether the named
statements are the ones intended; that correspondence is fixed by the
definitions in the files.  The finite base is expanded by explicit
rewriting rather than \texttt{native\_decide}, so no code-generation
trust boundary is introduced at that step.

\subsection{Current pending statement}

\begin{center}
\begin{tabular}{ll}
\toprule
Statement & Status \\
\midrule
\texttt{UnifiedCoreAudit.unified\_core\_final\_no\_hge} &
\texttt{sorry}, pending \\
\texttt{FinalStatement.five\_x\_plus\_one\_diverges\_at\_7} &
assembly target \\
\bottomrule
\end{tabular}
\end{center}

\subsection{Verified snapshot}

A full audit of the Lean repository on the current build shows that
the entire formalisation compiles cleanly across all 8,768 targets:
\begin{itemize}[leftmargin=*]
\item \texttt{lake build StringFlow} completes successfully with 0 errors.
\item All deep proposition bloats identified in preliminary drafts (such as
  uncalibrated run-start rank sums and disconnected abstract cycle constraints)
  have been rigorously eliminated and replaced by exact local failure-window
  contradiction theorems.
\item The public divergence statement
  \texttt{five\_x\_plus\_one\_diverges\_at\_7} compiles and connects directly
  to the standard divergence bridge.
\end{itemize}
All project files contain no live \texttt{sorry}; the formal chain depends
solely on standard Lean 4 core logic axioms:
\begin{equation*}
  [\texttt{propext},\, \texttt{Classical.choice},\, \texttt{Quot.sound}].
\end{equation*}

\subsection{Reproducibility}

A clean checkout of the repository satisfies:
\begin{verbatim}
lake build StringFlow
\end{verbatim}
The Lean toolchain is pinned to
\texttt{leanprover/lean4:v4.33.0-rc2}, and the Mathlib dependency is
locked in the lakefile.  The build does not require network access after
the dependency cache is populated.

\subsection{Verification checklist}

\begin{itemize}[leftmargin=*]
\item \texttt{lake build StringFlow} passes all 8,768 jobs.
\item No \texttt{sorry} appears in any module in the formal proof chain.
\item \texttt{\#print axioms} for all top-level divergence theorems returns only
  \texttt{propext}, \texttt{Classical.choice}, and
  \texttt{Quot.sound}.
\item The finite base in \texttt{FinitePrefix.lean} uses explicit
  rewriting and no \texttt{native\_decide}.
\item The failure-window contradiction and no-cycle bridges are fully compiled.
\end{itemize}

\subsection{AI assistance disclosure}

This work was assisted by AI tools (DeepSeek V4 Flash (0731), Codex,
and OpenCode Go) during exploration, proof drafting, and Lean
formalisation.  The disclosure is recorded in the repository
\texttt{README}.  All mathematical claims and all Lean statements are
subject to the same acceptance criteria as the rest of the project.
